\chapter{Lương tri và biện minh}
\markboth{Lương tri và biện minh}{Conscience and Excuses}

\section{Một người biết mình sai nhưng cứ nói ai rồi cũng vậy.}
\begin{truyen}{Một người biết mình sai nhưng cứ nói ai rồi cũng vậy.}{Common does not mean right.}
\chuhoa{A}{nh tự trấn an bằng cách bảo chuyện đó ai cũng làm. Nhưng việc nhiều người cùng sai không biến cái sai thành bình thường đúng nghĩa. Lương tri yếu đi rất nhanh khi ta dùng số đông để ru ngủ mình.}
\textit{(He comforted himself by saying everyone does it. But many people doing wrong does not make it right. Conscience weakens quickly when we use the crowd to lull ourselves.)}
\end{truyen}
\begin{baihoc}
Sự phổ biến của một hành vi không phải giấy chứng nhận cho sự đúng đắn của nó.
\textit{(The popularity of a behavior is not a certificate of its correctness.)}
\end{baihoc}
\begin{ghinhoanh}\textbf{Short English to remember:}

Common does not mean right.
\end{ghinhoanh}
\begin{tuhocnhanh}
1. Em có đang lấy số đông làm cớ cho điều chưa đúng không? \textit{(Are you using the crowd as an excuse for something not right?)}

2. Điều gì trong em biết rõ đây là sai? \textit{(What inside you clearly knows this is wrong?)}
\end{tuhocnhanh}
\ngancach

\section{Một người bảo mình làm vậy chỉ vì hoàn cảnh.}
\begin{truyen}{Một người bảo mình làm vậy chỉ vì hoàn cảnh.}{Circumstance explains, but does not erase.}
\chuhoa{H}{oàn cảnh quả thật có thể đẩy con người tới những chỗ khó. Nhưng nếu ta biến nó thành lý do để khỏi chịu trách nhiệm, ta sẽ ở mãi trong sai lầm mà không trưởng thành được.}
\textit{(Circumstances can indeed push people into hard places. But if we use them to avoid responsibility, we remain inside the mistake and fail to mature.)}
\end{truyen}
\begin{baihoc}
Giải thích được vì sao mình sai không đồng nghĩa với việc mình không còn sai nữa.
\textit{(Being able to explain why you were wrong is not the same as no longer being wrong.)}
\end{baihoc}
\begin{ghinhoanh}\textbf{Short English to remember:}

Circumstance explains, but does not erase.
\end{ghinhoanh}
\begin{tuhocnhanh}
1. Em có đang giải thích quá nhiều thay vì sửa không? \textit{(Are you explaining too much instead of fixing things?)}

2. Bước chịu trách nhiệm nào em còn thiếu? \textit{(What step of responsibility are you still missing?)}
\end{tuhocnhanh}
\ngancach

\section{Một người trốn việc rồi tự nhủ mình cần nghỉ xứng đáng.}
\begin{truyen}{Một người trốn việc rồi tự nhủ mình cần nghỉ xứng đáng.}{Excuses wear moral perfume.}
\chuhoa{C}{ó những cái cớ nghe rất dịu: mình mệt, mình xứng đáng, mình chỉ lần này thôi. Nhưng càng dùng những lời thơm tho để che một thói quen xấu, ta càng khó nhìn thẳng vào nó.}
\textit{(Some excuses sound gentle: I am tired, I deserve it, just this once. But the more we perfume bad habits with pleasant words, the harder they are to face honestly.)}
\end{truyen}
\begin{baihoc}
Điều nguy hiểm của biện minh là nó làm cái sai nghe có vẻ tử tế hơn chính nó.
\textit{(The danger of excuses is that they make wrongdoing sound kinder than it is.)}
\end{baihoc}
\begin{ghinhoanh}\textbf{Short English to remember:}

Excuses often smell nicer than truth.
\end{ghinhoanh}
\begin{tuhocnhanh}
1. Em có đang dùng lời dễ nghe để che điều không tốt không? \textit{(Are you using gentle words to cover something unhealthy?)}

2. Nếu gọi đúng tên nó, nó là gì? \textit{(If you named it honestly, what would it be?)}
\end{tuhocnhanh}
\ngancach

\section{Một người làm tổn thương ai đó rồi bảo mình vô ý.}
\begin{truyen}{Một người làm tổn thương ai đó rồi bảo mình vô ý.}{Unintentional still has impact.}
\chuhoa{A}{nh không cố ý, điều đó có thể đúng. Nhưng người kia vẫn đau, và nỗi đau ấy không tự nhỏ đi chỉ vì ý định ban đầu không xấu. Trưởng thành là biết rằng vô ý không xóa trách nhiệm sửa chữa.}
\textit{(He did not mean to, and that may be true. But the other person was still hurt, and that pain does not shrink just because the original intention was not bad. Maturity means knowing that being unintentional does not erase responsibility.)}
\end{truyen}
\begin{baihoc}
Ý định tốt là quan trọng, nhưng hậu quả thật cũng quan trọng không kém.
\textit{(Good intentions matter, but real consequences matter no less.)}
\end{baihoc}
\begin{ghinhoanh}\textbf{Short English to remember:}

Unintentional still has impact.
\end{ghinhoanh}
\begin{tuhocnhanh}
1. Em có đang núp sau chữ vô ý không? \textit{(Are you hiding behind the phrase I did not mean it?)}

2. Em có thể sửa hậu quả đó bằng cách nào? \textit{(How can you repair that impact?)}
\end{tuhocnhanh}
\ngancach

\section{Một người biết việc mình làm trái lương tâm nhưng quen dần.}
\begin{truyen}{Một người biết việc mình làm trái lương tâm nhưng quen dần.}{Conscience grows quiet when ignored.}
\chuhoa{L}{ần đầu anh còn thấy cấn. Lần sau thì đỡ hơn. Dần dần điều từng khiến anh khó chịu trở thành bình thường. Không phải vì nó đúng hơn, mà vì lương tri đã bị bắt im quá nhiều lần.}
\textit{(The first time he felt uneasy. The next time less so. Over time what once disturbed him felt normal. Not because it became right, but because conscience had been silenced too many times.)}
\end{truyen}
\begin{baihoc}
Cái sai lặp lại nhiều không làm nó nhẹ đi; nó chỉ làm ta chai đi.
\textit{(Repeated wrongdoing does not become lighter; it only makes us numb.)}
\end{baihoc}
\begin{ghinhoanh}\textbf{Short English to remember:}

Conscience grows quiet when ignored.
\end{ghinhoanh}
\begin{tuhocnhanh}
1. Điều gì trong em từng thấy cấn nhưng giờ đã chai đi? \textit{(What once bothered your conscience but now feels numb?)}

2. Em cần dừng chỗ nào trước khi chai thêm? \textit{(Where do you need to stop before becoming even more numb?)}
\end{tuhocnhanh}
\ngancach